\section{Integración entre Blender y Python}
En este contexto, se define la integración entre Blender y Python como el mecanismo de comunicación e 
interoperabilidad entre el entorno externo de ejecución en Python (utilizado por el componente de extracción) 
y el intérprete de Python embebido dentro de la arquitectura de Blender.
\medskip

La vinculación entre ambos entornos se concreta mediante un protocolo de intercambio basado en la persistencia 
de datos en disco y una estructura jerárquica de directorios. Este procedimiento opera bajo la siguiente lógica:

\begin{itemize}
    \item \textbf{Gestión de espacio de trabajo y proyectos:} El componente de extracción permite definir una 
    ruta de destino y un identificador único de proyecto, facilitando la organización, versión y diferenciación 
    entre distintas ejecuciones de la red.
    \item \textbf{Estructura de metadatos (JSON):} En la raíz del directorio del proyecto se genera un archivo de 
    configuración en formato \textit{JSON}. Este documento registra la topología de la arquitectura, la 
    secuencia de capas convolucionales, sus parámetros operativos y la ruta de acceso a los volúmenes 
    dispersos asociados, así como las métricas finales del modelo y otros metadatos.
    \item \textbf{Encapsulación de tensores (OpenVDB):} Los datos volumétricos derivados de los mapas de 
    activación y la salida del modelo se almacenan en archivos con formato \textit{OpenVDB}. La asociación 
    entre cada archivo de volumen y su correspondiente capa en la red se gestiona mediante una convención 
    estandarizada de nomenclatura.
\end{itemize}

El propósito de emplear el archivo \textit{JSON} es actuar como el manifiesto o archivo de configuración 
central del proyecto, concentrando la información estructural que Blender requiere para instanciar 
procedimentalmente la visualización. Al tratarse de un formato de texto plano y legible por humanos, 
facilita una auditoría rápida de la red previa a su renderizado. Asimismo, posibilita una visualización 
desacoplada y flexible, permitiendo generar esquemas de alto nivel de la topología sin necesidad de cargar 
la totalidad de los datos volumétricos en memoria.
\medskip

Esto ofrece la ventaja de permitir que el manifiesto actúe como un mecanismo de validación de datos. De forma 
previa a la generación gráfica, Blender emplea este archivo como el único punto de verdad 
(\textit{single source of truth}) para verificar la integridad estructural del proyecto. En caso de detectar 
anomalías, tales como archivos volumétricos faltantes o discrepancias en las dimensiones de los tensores, el 
manifiesto permite aplicar estrategias de resiliencia. En este escenario, el sistema puede sustituir el dato 
dañado por una representación geométrica abstracta alternativa, preservando la continuidad del flujo de trabajo 
y la visualización de la topología general sin interrumpir la ejecución del complemento.
\medskip

En secciones anteriores se detallaron las ventajas de emplear el formato OpenVDB para la encapsulación de los 
datos tensoriales extraídos en comparación con los formatos de imagen bidimensionales convencionales. 
Un beneficio técnico adicional reside en que, al tratarse de un estándar desarrollado originalmente para la 
industria de la animación y los efectos visuales, OpenVDB integra soporte nativo para secuencias temporales 
de cuadrículas volumétricas dentro de un mismo contenedor o serie de archivos (\textit{VDB Cache Sequences}). 
Esta capacidad técnica amplía las fronteras del sistema, abriendo la posibilidad de visualizar la evolución 
dinámica del modelo a través del tiempo, como el proceso de ajuste de pesos durante el entrenamiento, o bien 
auditar la respuesta topológica de las capas al procesar secuencias de datos continuos o video de entrada.
